\chapter{Introduction}
\lhead[\fancyplain{}{\bfseries\thepage}]{\fancyplain{}{\bfseries\rightmark}}

Cryptography has long been regarded as a form of art, rooted in humanity's need to communicate while concealing information from unintended recipients. Historically, much of the research in the field was conducted in secrecy, mainly because cryptographic techniques could provide significant military or political advantages. As a result, developments in cryptography were often responses to events or activities that were either secret at the time or ceased to be so only later \cite{the-code-book}. The history of cryptography is a continuous battle between codemakers (cipher makers) and codebreakers (cipher breakers). The creation of strong cipher, as well as the decryption of existing ones, heavily relied on creativity, inspiration, and strokes of genius \cite{introduction-to-modern-cryptography}. There was little underlying theory to rely on, and no precise definitions for determining whether a cryptographic scheme was secure or insecure. These factors contributed to cryptography being regarded as an art for centuries. However, between the 1970s and 1980s, the spread of the internet and the availability of affordable digital hardware brought cryptography into everyday life. At the same time, theoretical advances \cite{new-directions-in-cryptography, a-method-for-obtaining-digital-signatures, probabilistic-encryption} in information theory and computer science laid the foundations for provably secure cryptosystems, transforming cryptography into a scientific discipline. Cryptography may thus be seen as a science derived from art, an appealing definition, but one that does not fully capture what cryptography has become today. Modern cryptography is a field at the intersection of mathematics and computer science, concerned with techniques for ensuring data integrity, exchanging secret keys, authenticating users, enabling electronic auctions and elections, performing secure distributed computations, implementing digital cash, establishing server identities, and much more \cite{introduction-to-modern-cryptography}. While cryptography provides concrete solutions to these problems, such solutions are better viewed as applications of the field rather than the science itself. At its core, cryptography seeks solutions to clearly defined problems by abstracting away many of the complexities of real-world scenarios. This abstraction is necessary because cryptographic guarantees can only be made with respect to well-defined adversarial models and explicit assumptions. Although this approach may appear restrictive, the use of precise security definitions, rigorously specified assumptions, and mathematical proofs enables cryptography to deliver strong and reliable guarantees \cite{joyofcryptography}.

% \section{Security}
% A secrecy system is defined abstractly as a set of transformations of one
% space (the set of possible messages) into a second space (the set of possible
% cryptograms) \cite{}. Each particular transformation of the set corresponds to enciphering with a particular key. The transformations are supposed reversible
% (non-singular) so that unique deciphering is possible when the key is known.